% ID: SEL-D0054-OB-E02
% ESTADO: validada. S055.
\textbf{a)} Igualar ingresos y gastos equivale a $-3t^2+54t-96=0$, o $(t-2)(t-16)=0$. Coinciden en los años $\boxed{2\text{ y }16}$.

\textbf{b)} El beneficio, en miles de euros, es
\[
B(t)=I(t)-G(t)=-3t^2+54t-96=147-3(t-9)^2.
\]
La parábola abre hacia abajo. En $[0,18]$ corta el eje en $2$ y $16$ y toma el valor $-96$ en ambos extremos.
\begin{center}
\setlength{\unitlength}{0.9mm}
\begin{picture}(113,62)(-12,-4)
\qbezier(0,24)(49,24)(98,24)
\qbezier(0,0)(0,27)(0,54)
\qbezier(0,8)(45,89)(90,8)
\put(99,24){$t$}
\put(2,54){$B(t)$}
\put(8,18){\scriptsize $2$}
\put(78,18){\scriptsize $16$}
\put(40,51){\scriptsize $(9,147)$}
\put(90,2){\scriptsize $18$}
\end{picture}
\end{center}
\textbf{c)} El cuadrado se anula sólo en $t=9$, así que el beneficio máximo llega al cabo de $\boxed{9\text{ años}}$ y asciende a $\boxed{147\,000\text{ euros}}$.
